\section{Lecture II: Affine Springer fibers}\label{s:ASF}

Affine Springer fibers are analogs of Springer fibers for loop groups. They were introduced by Kazhdan and Lusztig \cite{KL}. Roughly speaking, in the case of classical groups, instead of classifying flags in a $k$-vector space fixed by a $k$-linear transformation, affine Springer fibers classify (chains of) lattices in an $F$-vector space fixed by an $F$-linear transformation, where $F=k\lr{t}$. The cohomology groups of affine Springer fibers carry actions of the affine Weyl group.

The setup in this section is the same as in \S\ref{ss:notation}.

\subsection{Loop group, parahoric subgroups and the affine flag variety}
Let $F=k\lr{t}$ be the field of formal Laurent series in one variable $t$.  Then $F$ has a discrete valuation $\val_{F}:F^{\times}\to\ZZ$ such that $\val_{F}(t)=1$ and its valuation ring is $\calO_{F}=k\tl{t}$.

\subsubsection{The loop group} The {\em positive loop group} $L^{+}G$ is a group-valued functor on $k$-algebras. For any $k$-algebra $R$, we define $LG(R):=G(R\tl{t})$. It turns out that $LG$ is representable by a scheme over $k$ which is not of finite type. 

For example, when $G=\GL_{n}$, an element in $LG(R)=\GL_{n}(R\tl{t})$ is given by $n^{2}$ formal Laurent series $a_{ij}=\sum_{s\geq 0}a^{(s)}_{ij}t^{s}$ ($1\leq i,j\leq n$), with $a^{(s)}_{ij}\in R$, subject to one condition that $\det((a_{ij}))$ (which is a polynomial in the $a^{(s)}_{ij}$ of degree $n$) is invertible in $R$. Therefore in this case $LG$ is an open subscheme in the infinite dimensional affine space with coordinates $a^{(s)}_{ij}$, $1\leq i,j\leq n$ and $s\geq 0$. 

Similarly we may define the {\em loop group} $LG$ to be the functor $LG(R)=G(R\lr{t})$ on $k$-algebras $R$. The functor $LG$ is no longer representable by a scheme, but rather by an {\em ind-scheme}. An ind-scheme is an inductive limit $\varinjlim_{m} X_{m}$ in the category of schemes, i.e., $\{X_{m}\}$ form an inductive system of schemes over $k$, and $\varinjlim_{m} X_{m}$ is the functor $R\mapsto \varinjlim_{m} X_{m}(R)$.  When $G=\GL_{n}$, we may  define $X_{m}$ to be the subfunctor of $LG$ such that $X_{m}(R)$ consists of $n$-by-$n$ invertible matrices with entries in $t^{-m}R\tl{t}\subset R\lr{t}$. Then the same argument as in the  $L^{+}G$ case shows that $X_{m}$ is representable by a scheme over $k$. For $m<m'$, we have a natural closed embedding $X_{m}\incl X_{m'}$, and $LG$ in this case is the inductive limit $\varinjlim_{m}X_{m}$.  For general $G$, see \cite[\S1]{BL} and \cite[\S2]{Faltings}.


\subsubsection{The affine Grassmannian}\label{sss:Gr} The affine Grassmannian $\Gr_{G}$ of $G$ is defined as the sheafification of the presheaf $R\mapsto LG(R)/L^{+}G(R)$ in the category of $k$-algebras under the fpqc topology.
In particular, we have $\Gr_{G}(k)=G(F)/G(\calO_{F})$. 

When $G=\GL_{n}$, the affine Grassmannian $\Gr_{G}$ can be identified with the moduli space of projective $R\tl{t}$-submodules $\L\subset R\lr{t}^{n}$ such that
\begin{equation}\label{lattice range}
(t^{m}R\tl{t})^{n}\subset \L\subset (t^{-m}R\tl{t})^{n}
\end{equation}
for some $m\geq 0$. Such an $R\tl{t}$-module $\L$ is called a {\em lattice} in $R\lr{t}^{n}$.  For fixed $m$, let $X_{m}$ be the subfunctor of $\Gr_{G}$ classifying those $\L$ such that \eqref{lattice range} holds, then $X_{m}$ is representable by a projective scheme over $k$. The natural closed embeddings $X_{m}\incl X_{m+1}$ make $\{X_{m}\}$ into an inductive system of projective schemes, and $\Gr_{G}$ is representable by the ind-scheme $\varinjlim_{m}X_{m}$.

Let us elaborate on the bijection between $\Gr_{G}(k)$ and the lattices in $F^{n}=k\lr{t}^{n}$. Let $\calO_{F}^{n}\subset F^{n}$ be the standard lattice. Let $\frL_{n}$ be the set of lattices in $F^{n}$ (in the case $R=k$ a lattice is simply an $\calO_{F}$-submodules of $F^{n}$ of rank $n$). The group $LG(k)=G(F)$ acts on $\frL_{n}$ by $LG\ni g:\L\mapsto g\L$. This action is transitive and the stabilizer of the standard lattice $\calO^{n}_{F}$ is $L^{+}G(k)=G(\calO_{F})$. Therefore this action induces a $G(F)$-equivariant bijection
\begin{equation}\label{Gr lattices}
\Gr_{G}(k)=G(F)/G(\calO_{F})\isom \frL_{n}.
\end{equation}


For general $G$, $\Gr_{G}$ is always representable by an ind-scheme of the form $\varinjlim_{m}X_{m}$ where $X_{m}$ are projective schemes over $k$, and the transition maps $X_{m}\incl X_{m+1}$ are closed embeddings.  We have a canonical exhaustion of $\Gr_{G}$ by projective schemes given by the affine Schubert stratification, which we now recall. The action of $L^{+}G$ on $\Gr_{G}$ by left translation has orbits indexed by dominant cocharacters $\l\in\xcoch(T)^{+}$. We denote by $\Gr_{G,\l}$ the $L^{+}G$-orbit through $\l(t)$. Let $\Gr_{G,\leq \l}$ be the closure of $\Gr_{G,\l}$. Then $\Gr_{G,\leq\l}$ is a projective scheme and $\Gr_{G}$ is the union of $\Gr_{G,\leq\l}$. For more details on the affine Grassmannian, we refer to \cite[\S2]{BL}, \cite[\S2]{Faltings} and Zhu's lectures \cite{Zhu}.




\subsubsection{Parahoric subgroups} The subgroup $L^{+}G$ of $LG$ is an example of a class of subgroups called {\em parahoric subgroups}.  Fix a Borel subgroup $B\subset G$ and let $\bI\subset L^{+}G$ be the preimage of $B$ under the map $L^{+}G\to G$ given by reduction modulo $t$. Then $\bI$ is an example of an {\em Iwahori subgroup} of $LG$. General Iwahori subgroups are conjugates of $\bI$ in $LG$. Like $L^{+}G$, Iwahori subgroups are group subschemes of $LG$ of infinite type. Parahoric subgroups are connected group subschemes of $LG$ containing an Iwahori subgroup with finite codimension. A precise definition of parahoric subgroups involves a fair amount of Bruhat-Tits theory, which we refer the readers to the original papers of Bruhat and Tits \cite{BT1}, and the survey paper \cite{Tits}. 

Just as the conjugacy classes of parabolic subgroups of $G$ are in bijection with subsets of the Dynkin diagram of $G$, the $LG$-conjugacy classes of parahoric subgroups of $LG$ are in bijection with {\em proper subsets} of the vertices of the {\em extended Dynkin diagram} $\wt{\Dyn}(G)$ of $G$, which has one more vertex than the Dynkin diagram of $G$. See Kac's book \cite[\S4.8]{Kac}, Bourbaki \cite[Ch VI]{Bourbaki} for extended Dynkin diagrams and the expository paper of Gross \cite{Gross} for connection with parahoric subgroups. 

Each $\bP$ admits a canonical exact sequence of group schemes
\begin{equation*}
1\to \bP^{+}\to \bP\to L_{\bP}\to1
\end{equation*}
where $\bP^{+}$ is the {\em pro-unipotent radical} of $\bP$ and $L_{\bP}$ is a reductive group over $k$, called the {\em Levi quotient} of $\bP$. If $\bP$ corresponds to a subset $J$ of the vertices of $\wt{\Dyn}(G)$, then the Dynkin diagram of the Levi quotient $L_{\bP}$ is the sub-diagram of $\wt{\Dyn}(G)$ spanned by $J$.





\subsubsection{Affine flag varieties} For each parahoric subgroup $\bP\subset LG$ we may define the corresponding {\em affine partial flag variety} $\Fl_{\bP}$ as the fpqc sheafification of the functor $R\mapsto LG(R)/\bP(R)$ on the category of $k$-algebras. This functor is also representable by an ind-scheme $\varinjlim_{m}X_{m}$ where each $X_{m}$ is a projective scheme over $k$ and the transition maps are closed embeddings.  The affine Grassmannian $\Gr_{G}$ is a special case of $\Fl_{\bP}$ for $\bP=L^{+}G$.

Consider the special case $\bP=\bI$ is an Iwahori subgroup of $LG$. When $G$ is simply-connected, $\bI$ is its own normalizer, and we may identify $\Fl_{\bI}$ as the moduli space of Iwahori subgroups of $LG$, hence giving an intrinsic definition of the affine flag variety. We usually denote $\Fl_{\bI}$ by $\Fl$ or $\Fl_{G}$ and call it {\em the affine flag variety of $G$}, with the caveat that $\Fl_{\bI}$ is canonically independent of the choice of $\bI$ only when $G$ is simply-connected.

Let $\bP\subset\bQ$ be two parahoric subgroups of $LG$. Then we have a natural projection $\Fl_{\bP}\to \Fl_{\bQ}$. The fibers of this projection are isomorphic to the partial flag variety of $L_{\bQ}$ corresponding to its parabolic subgroup given by the image of $\bP\to L_{\bQ}$. In particular, there is a natural projection $\Fl_{G}\to \Gr_{G}$ whose fibers are isomorphic to the flag variety $\sB$.


\subsubsection{The case of $\SL_{n}$}\label{sss:GLn parahoric} We have seen in \S\ref{sss:Gr} that the affine Grassmannian of $\GL_{n}$ has an interpretation as the moduli space of lattices. In fact, parahoric subgroups of $LG$ and the associated affine partial flag varieties can also be described using lattices. Here we consider the case $G=\SL_{n}$.  

Recall that the set of lattices in $F^{n}$ is denoted by $\frL_{n}$. For any two lattices $\L_{1},\L_{2}\in\frL_{n}$ we may define their relative length to be the integer
\begin{equation*}
[\L_{1}:\L_{2}]:=\dim_{k}(\L_{1}/\L_{1}\cap\L_{2})-\dim_{k}(\L_{2}/\L_{1}\cap\L_{2}).
\end{equation*}
Let $J\subset \ZZ/n\ZZ$ be a non-empty subset. Let $\tilJ$ be the preimage of $J$ under the projection  $\ZZ\to\ZZ/n\ZZ$. A {\em periodic $J$-chain of lattices} is a function $\L:\tilJ\to \frL_{n}$ sending each $i\in \tilJ$ to a lattice $\L_{i}\in\frL_{n}$ such that
\begin{itemize}
\item $[\L_{i}:\calO_{F}^{n}]=i$  for all $i\in \tilJ$;
\item $\L_{i}\subset \L_{j}$ for $i<j$ in $\tilJ$;
\item $\L_{i}=t\L_{i+n}$ for all $i\in \tilJ$.
\end{itemize}
Let $\frL_{J}$ be the set of periodic $J$-chains of lattices. For each $\{\L_{i}\}_{i\in \tilJ}\in\frL_{J}$, let $\bP_{\{\L_{i}\}_{i\in J}}\subset LG$ to be the simultaneous stabilizers of all $\L_{i}$'s. Then $\bP_{\{\L_{i}\}_{i\in J}}$ is a parahoric subgroup of $LG$. We call such a parahoric subgroup of type $J$. In fact all parahoric subgroups of $LG$ arise from a unique periodic $J$-chain of lattices, for a unique non-empty $J\subset \ZZ/n\ZZ$.  Therefore we get a bijection between $\sqcup_{J}\frL_{J}$ and the set of parahoric subgroups of $LG$.  In particular, $L^{+}G$ is the parahoric subgroup corresponding to the periodic $\{0\}$-chain of lattices given by $\L_{i}=t^{i/n}\calO^{n}_{F}$, where $i\in \tilJ=n\ZZ$. 

The extended Dynkin diagram of $G$ is a loop with $n$ nodes which we index cyclically by the set $\ZZ/n\ZZ$, such that $0$ corresponds to the extra node compared to the usual Dynkin diagram. Parahoric subgroups of type $J\neq\varnothing$ corresponds to the proper subset $\ZZ/n\ZZ-J$ of the nodes of the extended Dynkin diagram.
 
One can find the moduli space $\Fl_{J}$ of periodic $J$-chains of lattices such that  $\Fl_{J}(k)=\frL_{J}$. Fixing any parahoric subgroup $\bP$ of type $J$, $\Fl_{J}$ can be identified with the affine partial flag variety $\Fl_{\bP}$.  In particular, the affine flag variety $\Fl$ for $G=\SL_{n}$ can be identified with the moduli space of {\em periodic full chains of lattices}, i.e., a sequence of lattices $\cdots \L_{-1}\subset \L_{0}\subset \L_{1}\cdots$ in $F^{n}$ with $[\L_{i}:\calO^{n}_{F}]=i$ and $\L_{i}=t\L_{i+n}$ for all $i\in\ZZ$.



\subsubsection{The case of $\Sp_{2n}$}\label{sss:Sp parahoric} Now consider $G=\Sp_{2n}=\Sp(V)$, where $V=k^{2n}$ is equipped with a symplectic form. We extended the symplectic form on $V$ $F$-linearly to a symplectic form $\jiao{-,-}$ on $V\otimes_{k}F$. For a lattice $\L\in \frL_{2n}$,  define its symplectic dual to be the set $\L^{\vee}:=\{v\in V\otimes_{k}F|\jiao{v,\L}\subset \calO_{F}\}$. Then $\L^{\vee}$ is again a lattice in $V\otimes_{k}F$. The operation $\L\mapsto \L^{\vee}$ defines an involution on $\frL_{2n}$.  

Let $J\subset\ZZ/2n\ZZ$ be a non-empty subset stable under multiplication by $-1$. Let $\tilJ\subset \ZZ$ be the preimage of $J$ under the natural projection $\ZZ\to\ZZ/2n\ZZ$. A {\em periodic self-dual $J$-chain of lattices} in $V\otimes_{k}F$ is a periodic $J$-chain of lattices (i.e., an element in $\frL_{J}$ in the notation of \S\ref{sss:GLn parahoric}) satisfying the extra condition that
\begin{equation*}
\L^{\vee}_{i}=\L_{-i}, \textup{ for all } i\in \tilJ.
\end{equation*}
Denote the set of periodic self-dual $J$-chains of lattices in $V\otimes_{k}F$ by $\frL^{\Sp(V)}_{J}$. This is a set with an action of $G(F)=\Sp(V\otimes_{k}F)$. For any $\{\L_{i}\}_{i\in\tilJ}\in \frL^{\Sp(V)}_{J}$, the simultaneous stabilizer of the $\L_{i}$'s is a parahoric subgroup of $LG$, and every parahoric subgroup of $LG$ arises this way. For a parahoric subgroup $\bP$ of type $J$, the corresponding affine partial flag variety $\Fl_{\bP}$ can be identified with the moduli space of periodic self-dual $J$-chains of lattices so that $\Fl_{\bP}(k)\cong\frL^{\Sp(V)}_{J}$ as $G(F)$-sets. The readers are invited to work out the similar story for orthogonal groups, see Exercise \ref{exII:SO parahoric}.





